% Generated by scripts/export_feature_tables.py -- do not edit by hand.
\begin{description}
  \item[\texttt{user\_activity\_count}] Interaction volume separates a user with enough history for collaborative signal from one riding on the popularity prior. The distribution is heavy-tailed, so it is log-compressed before any model consumes it linearly.

  \item[\texttt{user\_mean\_rating}] Implicit filtering keeps only positives, so this measures how generously a user rates things they already liked. It calibrates the confidence attached to a positive rather than the positive itself.

  \item[\texttt{user\_tenure\_days}] Calendar span between a user's first and last training event. Distinguishes fifty ratings spread over three years from fifty ratings in one weekend; the two imply very different taste stability.

  \item[\texttt{user\_recency\_days}] Days from a user's last training event to the cut-off. A dormant user's stored preferences are staler than an active user's, and the serving path should weigh them accordingly. Depends on the cut-off, so it is recomputed per partition and never cached across splits.

  \item[\texttt{user\_genre\_affinity}] Genre preferences weighted by 2^(-age\_days / half\_life), so a genre abandoned two years ago is discounted against one watched last month. MovieLens has true event times, which makes the decay calendar-based rather than a rank proxy.

  \item[\texttt{item\_popularity\_rank}] Normalized dense popularity rank over the full catalogue. Backs the popularity baseline and the cold-start fallback. Total over the catalogue by construction, so serving never hits a missing key.

  \item[\texttt{item\_trailing\_popularity}] Interaction count in the 90 days before the cut-off. Lifetime popularity is dominated by whatever was popular years ago; a trailing window tracks what is popular now, which is what a system serving traffic at the cut-off actually needs. This is the feature that justifies evaluating on a temporal split.

  \item[\texttt{item\_genre\_tokens}] Sorted, deduplicated genre tokens. Content signal available for items with no interactions at all, which is the only thing a purely collaborative model cannot supply for a cold item.

  \item[\texttt{item\_release\_year}] Release year parsed from the trailing "(YYYY)" in the MovieLens title, which has no dedicated column. Combined with the event timestamp it yields item age at interaction time, a strong driver of catalogue demand. The fallback is a sentinel meaning "unknown", not year zero.

\end{description}
